% ============================================================================
\section{Example IV: the ammonia inversion doublet}
\label{sec:ammonia}
% ============================================================================

The three examples so far built their two-level space from two places an electron can
be. This one builds it from two \emph{shapes the same molecule can have}.

\subsection{Tunneling as an off-diagonal matrix element}

Ammonia is pyramidal, and the nitrogen can pass through the plane of the three hydrogens
like an umbrella turning inside out. Label the two pyramidal geometries $\ket{L}$
(umbrella down) and $\ket{R}$ (umbrella up). Classically they are separated by a barrier
of roughly \SI{2000}{\wavenumber} and the molecule would be stuck in one of them.
Quantum mechanically the nitrogen \textbf{tunnels}, and tunneling is nothing other than
an off-diagonal matrix element:
%
\begin{equation}
\Hop = \begin{bmatrix} 0 & -\Delta_t \\ -\Delta_t & 0\end{bmatrix} = -\Delta_t\,\sxo .
\label{eq:nh3-bare}
\end{equation}
%
Pure $\sxo$. The eigenstates are $\tfrac{1}{\sqrt2}(\ket{L}\pm\ket{R})$, split by
$2\Delta_t$: the \textbf{inversion doublet}, measured at
%
\begin{equation}
2\Delta_t = \SI{0.7934}{\wavenumber} = \SI{23.786}{\giga\hertz},
\end{equation}
%
a microwave transition --- the line Gordon, Zeiger and Townes used to build the first
\textbf{maser} in 1954, three years before the laser \cite{Townes1955}.

The conceptual point is worth dwelling on. The energy eigenstates of ammonia are
\emph{not} the pyramidal geometries; they are delocalized over both. A molecule prepared
as $\ket{L}$ is a superposition of the two eigenstates, so by \cref{sec:bloch} it
oscillates: the umbrella flips back and forth at \SI{23.786}{\giga\hertz}.

\subsection{Adding a $\szo$ term: the Stark effect}

Switch on a static electric field $\mathcal{E}$ along the molecular axis. The two
pyramidal geometries have dipole moments pointing in \emph{opposite} directions, so the
field stabilizes one and destabilizes the other:

\begin{keyresult}[Ammonia in a field]
\vspace{-0.35cm}
\begin{equation}
\Hop = \mu\mathcal{E}\,\szo - \Delta_t\,\sxo
\qquad\Longrightarrow\qquad
d_z = 2\mu\mathcal{E}, \quad d_x = -2\Delta_t ,
\label{eq:nh3-stark}
\end{equation}
with $\mu = \SI{1.47}{\debye}$ the NH$_3$ dipole moment.
\end{keyresult}

\noindent
Structurally this is \cref{eq:dimer-d} again --- but now \emph{you} control $d_z$ with a
knob on the bench. The gap is
%
\begin{equation}
E_+-E_- = \sqrt{(2\mu\mathcal{E})^2+(2\Delta_t)^2},
\end{equation}
%
quadratic in $\mathcal{E}$ at low field (a second-order Stark effect) and crossing over
to linear at high field. The crossover occurs where the two terms are comparable,
$\mu\mathcal{E}=\Delta_t$, that is at
%
\begin{equation}
\mathcal{E}_{\rm cross} = \frac{\Delta_t}{\mu} = \SI{16.07}{\kilo\volt\per\centi\meter},
\end{equation}
%
at which point the gap has grown by exactly $\sqrt{2}$. Turn the field well past this
and $\theta\to0$: the field literally traps the umbrella on one side, converting a
delocalized molecule into a localized one. \Cref{fig:nh3} shows both halves.

% ---------------------------------------------------------------------------
\begin{figure}[t]
\centering
\begin{tikzpicture}
\begin{groupplot}[
  group style={group size=2 by 1, horizontal sep=1.6cm},
  notestyle, height=6.0cm,
  xmin=0, xmax=60,
]
\nextgroupplot[
  xlabel={electric field $\mathcal{E}$ (kV\,cm$^{-1}$)},
  ylabel={energy (cm$^{-1}$)},
  title={Stark effect on the doublet},
  ymin=-1.7, ymax=1.7,
  legend style={at={(0.03,0.97)},anchor=north west,draw=none,fill=white,
                fill opacity=0.85,text opacity=1,font=\footnotesize},
  legend cell align=left,
]
% mu = 0.024685 cm^-1 per kV/cm ; Delta_t = 0.3967 cm^-1
\addplot[cUpper, very thick, domain=0:60, samples=200]
  {sqrt((0.024685*x)^2 + 0.3967^2)};
\addlegendentry{upper level}
\addplot[cLower, very thick, domain=0:60, samples=200]
  {-sqrt((0.024685*x)^2 + 0.3967^2)};
\addlegendentry{lower level}
\addplot[cMuted, dashed, thin, domain=0:60, samples=2]{0.024685*x};
\addlegendentry{$\pm\mu\mathcal{E}$ (no tunneling)}
\addplot[cMuted, dashed, thin, domain=0:60, samples=2]{-0.024685*x};
\draw[cMuted,densely dotted,thin] (axis cs:16.07,-1.7) -- (axis cs:16.07,1.7);
\draw[<->,cAccent,thick] (axis cs:2.6,-0.3967) -- (axis cs:2.6,0.3967);
\node[cAccent,font=\footnotesize,anchor=west] at (axis cs:3.2,0.72){$2\Delta_t$};

\nextgroupplot[
  xlabel={electric field $\mathcal{E}$ (kV\,cm$^{-1}$)},
  ylabel={$|\Braket{\szo}|$ of the lower state},
  title={The field traps the umbrella},
  ymin=-0.03, ymax=1.03,
]
\addplot[cAccent, very thick, domain=0:60, samples=300]
  {abs(0.024685*x)/sqrt((0.024685*x)^2 + 0.3967^2)};
\draw[cMuted,densely dotted,thin] (axis cs:16.07,-0.03) -- (axis cs:16.07,1.03);
\node[cMuted,font=\footnotesize,anchor=west] at (axis cs:17.5,0.20)
  {$\mu\mathcal{E}=\Delta_t$};
\node[font=\footnotesize,anchor=west,align=left] at (axis cs:1.5,0.86)
  {delocalized $\rightarrow$ localized};
\end{groupplot}
\end{tikzpicture}
\caption{Ammonia in a static electric field, from \cref{eq:nh3-stark}. \textbf{Left:}
the levels start split by the tunneling doublet $2\Delta_t = \SI{0.7934}{\wavenumber}$
and approach the classical Stark asymptotes $\pm\mu\mathcal{E}$ at high field.
\textbf{Right:} the localization $|\Braket{\szo}| = |d_z|/|\dvec| = |\cos\theta|$ of the
lower state, running from $0$ (equal amplitude on both umbrella geometries) to $1$
(trapped on one side).}
\label{fig:nh3}
\end{figure}
% ---------------------------------------------------------------------------

\begin{nbbox}[A puzzle worth posing to yourself]
At zero field the ground state is $\tfrac{1}{\sqrt2}(\ket{L}+\ket{R})$, which has
\emph{zero} dipole moment --- yet NH$_3$ is a famously polar molecule with
$\mu=\SI{1.47}{\debye}$. Both statements are true. The resolution turns on what a
dipole-moment measurement actually averages over, and on what timescale; see
Problem~\ref{prob:nh3} in \cref{sec:problems}.
\end{nbbox}
